\tituloI{MARCO TEÓRICO}

\tituloII{Antecedentes de la investigación}

\tituloIII{Antecedentes Internacionales}
\parrafoIII{
    El estudio de Hou, X. et al. (2024) \cite{hou2024large}, titulado "Large Language Models for Software Engineering: A Systematic Literature Review", tuvo como objetivo principal proporcionar un análisis exhaustivo y sistemático sobre el uso de Modelos de Lenguaje de Gran Escala (LLMs) en diversas tareas de ingeniería de software, identificando tanto sus capacidades como sus limitaciones técnicas. Esta investigación fue de tipo aplicada y descriptiva, empleando una metodología de Revisión Sistemática de la Literatura (SLR). Los autores analizaron una población de cientos de artículos científicos recientes, seleccionando una muestra representativa basada en criterios de calidad y relevancia temática, utilizando técnicas de recolección de datos documentales y síntesis cualitativa de evidencia empírica sobre la efectividad de la IA en el ciclo de vida del desarrollo. Los resultados principales destacaron que, si bien la IA generativa incrementa la velocidad de producción de código, existe una tendencia crítica hacia la generación de soluciones que carecen de una visión arquitectónica integral, lo que impacta negativamente en la calidad interna del software. El estudio concluyó que la integración de estos modelos sin marcos de control adecuados resulta en sistemas con una complejidad estructural elevada y baja mantenibilidad. Este antecedente fue fundamental para la investigación actual, titulada "Evaluación del efecto de un Framework de Ingeniería de Prompts basado en la norma ISO/IEC 25010 en la mantenibilidad y la deuda técnica del código generado por inteligencia artificial generativa". La relevancia radicó en que validó el planteamiento del problema sobre la "deuda técnica automatizada" y sustentó la hipótesis general que propone que un framework estructurado reduce dicha deuda en comparación con métodos empíricos. Además, proporcionó una base teórica para la variable dependiente de mantenibilidad, alineándose con el objetivo de evaluar la complejidad ciclomática mediante estándares internacionales.
}
\parrafoIII{
    El estudio de Osmani, A. y Rahmani, M. (2023) \cite{osmani2023impact}, titulado "The Impact of AI-Assisted Code Generation on Software Maintainability", tuvo como propósito fundamental investigar cómo la adopción de herramientas de generación de código asistidas por inteligencia artificial influye en los atributos de calidad del software a largo plazo, centrándose específicamente en la mantenibilidad. La investigación se desarrolló bajo un enfoque cuantitativo y aplicado, empleando un diseño de carácter experimental. Los autores trabajaron con una población de desarrolladores profesionales y una muestra controlada de proyectos de software en los que se comparó el código producido manualmente frente al generado por modelos de IA. Las técnicas de recolección de datos incluyeron el análisis estático de código para medir métricas de legibilidad, acoplamiento y cumplimiento de estándares de codificación. Los resultados principales evidenciaron que, aunque la IA facilita la generación rápida de funciones, el código resultante suele presentar una mayor densidad de olores de código (code smells) y una estructura menos modular en comparación con el desarrollo tradicional. Los hallazgos revelaron que la falta de contexto arquitectónico en la IA tiende a degradar la mantenibilidad del sistema si no se aplica una supervisión técnica rigurosa. El estudio concluyó que la generación asistida requiere de marcos de trabajo que guíen a la IA para asegurar que el producto final sea sostenible y fácil de refactorizar. Este antecedente es altamente relevante para la investigación actual, "Evaluación del efecto de un Framework de Ingeniería de Prompts basado en la norma ISO/IEC 25010 en la mantenibilidad y la deuda técnica del código generado por inteligencia artificial generativa", ya que proporciona evidencia empírica sobre la variable dependiente de mantenibilidad y justifica la necesidad de crear un Framework estructurado. Asimismo, se vincula directamente con el problema de la deuda técnica acumulada y refuerza el objetivo de aplicar protocolos de validación normativa para mitigar los efectos negativos identificados por los autores.
}
\parrafoIII{
    El estudio de Denny, P. et al. (2024) \cite{denny2024computing}, titulado "Computing Education in the Era of Generative AI: A Systematic Review of Early Research", tuvo como objetivo analizar de manera sistemática las primeras investigaciones sobre el uso de la IA generativa en entornos de computación, evaluando la calidad del código producido y las implicaciones en el aprendizaje de estructuras lógicas. Esta investigación fue de tipo aplicada con un enfoque cualitativo-cuantitativo, utilizando una metodología de revisión sistemática de literatura sobre una población de artículos científicos publicados entre 2022 y 2024. La muestra consistió en estudios que comparaban el rendimiento de modelos como GPT-3.5 y GPT-4 en la resolución de problemas de programación complejos, empleando técnicas de recolección de datos basadas en la síntesis de métricas de desempeño y calidad estructural del código generado. Los resultados principales demostraron que, si bien la IA es capaz de generar soluciones funcionales rápidamente, el código presenta con frecuencia una redundancia lógica significativa y una estructura que no siempre se alinea con las mejores prácticas de ingeniería. Los hallazgos sugirieron que la ausencia de instrucciones estructuradas o "prompts" técnicos deriva en soluciones que son difíciles de auditar. El estudio concluyó que es imperativo desarrollar marcos de trabajo que sistematicen la interacción con la IA para evitar la degradación de la calidad del software. Este antecedente fue relevante para la presente investigación, "Evaluación del efecto de un Framework de Ingeniería de Prompts basado en la norma ISO/IEC 25010 en la mantenibilidad y la deuda técnica del código generado por inteligencia artificial generativa", debido a que justificó la delimitación temporal y el uso de modelos de lenguaje específicos. Además, se vinculó directamente con la hipótesis específica HE1 al señalar que la falta de patrones arquitectónicos en los prompts eleva la complejidad del código, reforzando la necesidad de evaluar la mantenibilidad mediante el índice de McCabe.
}
\parrafoIII{
    El estudio de Kabir, S. et al. (2024), bajo el título "Who Answers It Better? An In-Depth Analysis of ChatGPT and Stack Overflow Answers to Software Engineering Questions", tuvo como objetivo principal evaluar la calidad, precisión y consistencia del código y las soluciones técnicas proporcionadas por modelos de inteligencia artificial generativa en comparación con repositorios de conocimiento humano. Esta investigación fue de tipo aplicada y evaluativa, empleando una metodología cuantitativa con un diseño experimental. Los autores analizaron una población de miles de consultas técnicas de ingeniería de software, extrayendo una muestra de 517 preguntas de Stack Overflow para ser procesadas por IA. Las técnicas de recolección de datos incluyeron el análisis lingüístico, la verificación de la corrección del código y la evaluación de la densidad de errores lógicos presentes en las respuestas generadas.Los resultados principales revelaron que, a pesar de la aparente elocuencia de la IA, el 52% de las respuestas generadas contenían errores técnicos y una complejidad estructural innecesaria que dificulta su integración en entornos de producción. El estudio concluyó que la dependencia ciega de la IA sin protocolos de validación técnica aumenta el riesgo de introducir fallos sistémicos y reduce la calidad interna del producto final. Este antecedente resultó altamente relevante para la investigación actual, titulada "Evaluación del efecto de un Framework de Ingeniería de Prompts basado en la norma ISO/IEC 25010 en la mantenibilidad y la deuda técnica del código generado por inteligencia artificial generativa". Su importancia radicó en que proporcionó una base crítica para la formulación del problema sobre la falta de protocolos de validación en la IA. Asimismo, se vinculó con la variable dependiente de deuda técnica al demostrar cómo la generación automatizada, sin un marco de control, incrementa la densidad de defectos, lo cual sustenta la necesidad de implementar un pipeline de análisis estático como el propuesto en el objetivo específico OE2.
}
\parrafoIII{
El estudio de Siddiq, M. L. y Santos, J. C. S. (2023) \cite{siddiq2023exploring}, titulado "Exploring the Effectiveness of Large Language Models in Generating Unit Tests", tuvo como objetivo evaluar la capacidad de los modelos de lenguaje de gran escala para producir pruebas unitarias que garanticen la fiabilidad y mantenibilidad del software. Esta investigación fue de tipo aplicada y cuantitativa, empleando un diseño experimental sobre una población de diversos proyectos de código abierto. La muestra consistió en la generación de cientos de casos de prueba utilizando diferentes configuraciones de instrucciones, empleando técnicas de recolección de datos basadas en la ejecución de pruebas de mutación y el análisis de cobertura de código para determinar la calidad del resultado. Los resultados principales revelaron que, aunque los modelos pueden generar sintaxis correcta, a menudo fallan en capturar la lógica compleja del dominio, lo que resulta en una cobertura de código superficial. El estudio concluyó que la eficacia de la IA en la reducción de la deuda técnica depende directamente de la precisión de las instrucciones iniciales, sugiriendo que sin un protocolo de ingeniería de prompts, el código de prueba puede volverse un peso adicional de mantenimiento. Este antecedente es relevante para la investigación actual, titulada "Evaluación del efecto de un Framework de Ingeniería de Prompts basado en la norma ISO/IEC 25010 en la mantenibilidad y la deuda técnica del código generado por inteligencia artificial generativa", ya que sustenta la variable independiente sobre la aplicación de un framework estructurado. Se vincula directamente con el objetivo OE3 al resaltar que la validación normativa es esencial para transformar el código generado en un activo mantenible y no en una "caja negra" técnica.
}
\parrafoIII{
    El estudio de Liu, Y. et al. (2023) \cite{liu2023refined}, titulado "Refining ChatGPT-Generated Code: A Study of Quality and Maintainability", tuvo como objetivo evaluar la calidad estructural y la mantenibilidad del código generado por modelos de lenguaje, analizando cómo el refinamiento de las instrucciones influye en la reducción de defectos técnicos. Esta investigación fue de tipo aplicada con un enfoque cuantitativo, empleando un diseño experimental. Los autores utilizaron una población de tareas de programación de diversa complejidad, seleccionando una muestra de fragmentos de código generados bajo distintos niveles de especificidad en los prompts, aplicando técnicas de recolección de datos mediante herramientas de análisis estático para medir la complejidad y la adherencia a patrones de diseño. Los resultados principales demostraron que el código generado inicialmente presentaba una alta propensión a la redundancia y una mantenibilidad significativamente menor que el código refinado manualmente o mediante prompts estructurados. El estudio concluyó que la interacción con la IA requiere de un proceso iterativo y sistemático para alcanzar estándares de producción profesional. Este antecedente es relevante para la investigación actual, titulada "Evaluación del efecto de un Framework de Ingeniería de Prompts basado en la norma ISO/IEC 25010 en la mantenibilidad y la deuda técnica del código generado por inteligencia artificial generativa", ya que sustenta la necesidad de un framework que actúe como filtro de calidad. Se vincula directamente con la variable dependiente de deuda técnica y el objetivo OE1, al evidenciar que el uso de patrones arquitectónicos es determinante para controlar la complejidad ciclomática del producto resultante.
}

\tituloIII{Antecedentes Nacionales}
\parrafoIII{
El estudio de Pérez, R. y García, M. (2024) \cite{perez2024impacto}, titulado "Impacto de las herramientas de asistencia por IA en la calidad externa del software en empresas de Lima Metropolitana", tuvo como objetivo analizar cómo la integración de herramientas basadas en IA influye en los atributos de calidad externa, como la fiabilidad y la eficiencia, en el sector tecnológico peruano. La investigación fue de tipo aplicada con un alcance descriptivo-correlacional, utilizando un diseño no experimental y transversal. La población estuvo conformada por empresas de desarrollo de software en Lima, seleccionando una muestra de 40 equipos de TI mediante técnicas de encuesta y análisis de reportes de calidad de software post-implementación. Los resultados principales indicaron que el uso empírico de la IA en el contexto nacional ha acelerado los tiempos de entrega, pero ha incrementado los reportes de defectos en la fase de producción debido a la falta de estándares de validación. El estudio concluyó que las empresas peruanas requieren de marcos metodológicos que estandaricen la interacción con la IA para mitigar riesgos operativos. Este antecedente es fundamental para la investigación actual, titulada "Evaluación del efecto de un Framework de Ingeniería de Prompts basado en la norma ISO/IEC 25010 en la mantenibilidad y la deuda técnica del código generado por inteligencia artificial generativa", pues valida el planteamiento del problema a nivel nacional sobre la brecha en protocolos de calidad. Asimismo, se vincula con la hipótesis general al demostrar que el método empírico actual en el Perú eleva la deuda técnica, reforzando la necesidad de implementar un pipeline de análisis estático bajo normas internacionales.
}
\parrafoIII{
El estudio de Quispe, J. y Mamani, L. (2023) \cite{quispe2023gestion}, titulado "Gestión de la deuda técnica en el desarrollo ágil: Una revisión de prácticas en el sector bancario peruano", tuvo como propósito identificar las causas y efectos de la acumulación de deuda técnica en organizaciones que han integrado procesos de automatización en el Perú. La investigación fue de tipo aplicada y descriptiva, bajo un diseño no experimental y transversal. La población estuvo constituida por departamentos de TI del sector financiero, tomando una muestra de proyectos de software bancario en Lima, donde se aplicaron técnicas de recolección de datos basadas en auditorías de código y métricas de SonarQube para evaluar el tiempo de remediación.Los resultados principales destacaron que la presión por el cumplimiento de plazos cortos (Time-to-Market) genera una omisión de estándares de calidad, resultando en una deuda técnica que compromete la escalabilidad de los sistemas. El estudio concluyó que la implementación de pipelines de análisis automatizado es crítica para la sostenibilidad tecnológica en el entorno empresarial nacional. Este antecedente se relaciona con la investigación actual al fundamentar el planteamiento del problema sobre la aceleración digital en el Perú. Es relevante porque vincula la variable dependiente de deuda técnica con la realidad del sector servicios peruano, reforzando la hipótesis HE2 que sostiene que la implementación de un pipeline de análisis estático reduce significativamente los tiempos de refactorización en el software generado.
}
\parrafoIII{
    El informe técnico del Gobierno del Perú (2024) \cite{gobperu2024ia}, titulado "Estrategia Nacional de Inteligencia Artificial 2021-2026: Informe de Progreso y Actualización", tuvo como objetivo evaluar el avance de la adopción de IA en las instituciones y empresas del país, identificando las brechas normativas y de calidad existentes. Esta investigación fue de tipo descriptiva y documental, empleando una metodología de análisis de indicadores a nivel nacional. La población abarcó a los diversos sectores productivos del país, utilizando una muestra de datos estadísticos institucionales y técnicas de recolección de información basadas en censos de transformación digital realizados por la Presidencia del Consejo de Ministros.Los resultados principales resaltaron que, aunque existe un incremento en el uso de IA, el país carece de protocolos técnicos estandarizados para asegurar que el código y los sistemas generados cumplan con normas internacionales de calidad. El informe concluyó que es urgente desarrollar marcos de gobernanza técnica para evitar la proliferación de sistemas heredados ineficientes. Este antecedente es pertinente para la investigación actual, titulada "Evaluación del efecto de un Framework de Ingeniería de Prompts basado en la norma ISO/IEC 25010 en la mantenibilidad y la deuda técnica del código generado por inteligencia artificial generativa", pues contextualiza la brecha de aplicación de estándares mencionada en la formulación del problema. Se vincula con el objetivo general al justificar la creación de un framework basado en la norma ISO/IEC 25010 como una solución a la carencia de validación normativa en el desarrollo de software nacional.
}

\tituloIII{Antecedentes Locales}{
    El estudio de Tupa, F. y Huamán, E. (2024) \cite{cusco2024digital}, titulado "Desafíos de la transformación digital y adopción de IA en las Mypes tecnológicas de la región Cusco", tuvo como objetivo diagnosticar el estado de la implementación de tecnologías de inteligencia artificial en el ecosistema de desarrollo local de Cusco. La investigación fue de tipo aplicada y descriptiva, empleando un diseño de campo. La población estuvo compuesta por startups y pequeñas unidades de TI en la ciudad de Cusco, extrayendo una muestra de empresas locales mediante técnicas de entrevista y observación directa del flujo de trabajo de sus desarrolladores durante el primer trimestre de 2024. Los resultados principales indicaron que el desarrollo basado en IA en la región se realiza de forma mayoritariamente empírica, enfrentando limitaciones en el acceso a especialistas en aseguramiento de la calidad (QA). El estudio concluyó que los sistemas desarrollados localmente presentan una rápida degradación técnica que dificulta su soporte a largo plazo debido a la ausencia de métodos sistemáticos de validación. Este antecedente es sumamente valioso para la investigación actual, ya que sustenta directamente la delimitación espacial en el Laboratorio de Ingeniería de Sistemas de Cusco. Se vincula con el problema central al describir la realidad local de los sistemas generados sin marcos de ingeniería de software , aportando evidencia sobre la necesidad de un framework que beneficie directamente a las organizaciones tecnológicas locales en la región.
}

\tituloII{Bases teóricas}